\section{Conclusion}

This paper presents a dynamic collaborative scheduling system for new energy logistics fleets based on graph algorithms and reinforcement learning. We model the urban road network as a weighted graph and build a discrete-time simulation framework that explicitly considers battery limits, load capacity, task deadlines, and charging station congestion. The system implements multiple path planning algorithms (Dijkstra, A*, RRT) and scheduling strategies, including fixed heuristics (nearest-task-first, maximum-weight-first, earliest-deadline-first) and a Q-learning-based hyper-heuristic that dynamically selects rules based on system state.

Our experimental evaluation across three problem scales reveals several key findings. First, among fixed heuristics, EDF and Nearest consistently outperform Heaviest, especially as task density and resource competition increase. Second, the Q-learning hyper-heuristic achieves competitive or superior performance compared to the best fixed baseline, with its main advantage lying in state-dependent dynamic rule switching rather than replacing a single heuristic. Third, training convergence analysis shows that Q-learning converges stably, though training difficulty increases with problem scale. The mixed-training model exhibits better cross-scale generalization despite sacrificing some peak performance. Fourth, charging strategy design significantly impacts overall system performance, with the optimal station strategy consistently outperforming the nearest station approach for both fixed heuristics and Q-learning methods.

The system follows a layered architecture that decouples road modeling, path planning, scheduling decisions, simulation execution, and visualization, enabling fair algorithm comparison and future extensions. Our work integrates graph algorithms, heuristic scheduling, reinforcement learning, and exact optimization into a unified framework, providing a practical platform for research on multi-agent coordination and complex urban delivery scenarios. The experimental results demonstrate that learning-based hyper-heuristics can effectively compete with or surpass fixed baselines through dynamic rule selection, and that charging strategy optimization must be jointly considered with task scheduling. As problem scale increases, the limitations of fixed rules become more evident, while learning-based methods show better adaptability, providing a foundation for future studies in new energy logistics scheduling.

